\newpage
\chapter{PENDAHULUAN} \label{Bab I}

\section{Latar Belakang} \label{I.LatarBelakang}
Bencana alam seperti gempa bumi, banjir, kebakaran hutan, dan badai
tropis menimbulkan dampak yang sangat merusak dan memerlukan respons
cepat dan terkoordinasi untuk menyelamatkan nyawa. Menurut
\textit{United Nations Office for the Coordination of Humanitarian
Affairs} (UNOCHA), 72 jam pertama setelah bencana merupakan periode
kritis dimana respons harus segera dimulai untuk memaksimalkan
penyelamatan korban \cite{unocha2017disaster}. Dalam periode tersebut,
tim respons darurat membutuhkan informasi akurat mengenai lokasi
korban, tingkat kerusakan infrastruktur, dan kebutuhan bantuan
kemanusiaan untuk mengalokasikan sumber daya secara efektif
\cite{braik2024automated}. \par

Media sosial telah menunjukkan potensi besar sebagai sumber informasi
kesadaran situasional selama kejadian bencana, dimana ribuan gambar dan
video dibagikan oleh masyarakat terdampak melalui platform seperti
\textit{Twitter} dan \textit{Facebook} dalam hitungan jam
\cite{seneviratne2024social}. Konten visual ini dapat memberikan
gambaran waktu nyata mengenai kondisi lapangan, membantu tim SAR
(\textit{Search and Rescue}) mengidentifikasi area prioritas, dan
memfasilitasi koordinasi bantuan kemanusiaan \cite{gopal2024leveraging}.
Namun, volume data yang sangat besar dan kebutuhan akan respons yang
cepat membuat kurasi manual tidak praktis untuk diterapkan, sehingga
diperlukan sistem klasifikasi otomatis berbasis \textit{machine
learning} untuk memfilter informasi yang relevan dari informasi yang
tidak berguna \cite{awasthi2023machine}. \par

Untuk memfasilitasi penelitian klasifikasi otomatis citra bencana,
Alam et al. memperkenalkan \textit{CrisisMMD}, sebuah \textit{dataset}
multimodal yang mengandung lebih dari 18.000 citra dari tujuh kejadian
bencana yang semuanya terjadi pada tahun 2017, yaitu \textit{Hurricane
Harvey}, \textit{Hurricane Irma}, \textit{Hurricane Maria}, gempa bumi
Mexico, kebakaran hutan California, gempa bumi Iraq-Iran, dan banjir
Sri Lanka \cite{alam2018crisismmd}. \textit{Dataset} ini menyediakan
tiga tugas klasifikasi utama: identifikasi keterdampakan informasi
(\textit{Informative Classification}), kategorisasi kemanusiaan
(\textit{Humanitarian Categories}), dan penilaian tingkat kerusakan
(\textit{Damage Severity Assessment}). \textit{CrisisMMD} telah menjadi
standar evaluasi dalam domain respons bencana dan digunakan secara luas
oleh komunitas peneliti untuk mengevaluasi berbagai pendekatan
komputasional \cite{firmansyah2022ensemble}. \par

Perkembangan \textit{deep learning} dalam \textit{computer vision}
telah melahirkan dua paradigma arsitektur yang berbeda secara mendasar
namun saling melengkapi. Di satu sisi, \textit{Convolutional Neural
Network} (CNN) telah menjadi paradigma dominan selama lebih dari satu
dekade berkat kemampuannya menangkap hierarki fitur visual secara lokal
\cite{lecun1998gradient}\cite{linardos2022deep}. Arsitektur CNN terus
berkembang, mulai dari pendekatan \textit{compound scaling} pada
keluarga \textit{EfficientNet} \cite{tan2019efficientnet} hingga
\textit{EfficientNetV2} yang menghadirkan perbaikan signifikan pada
kecepatan pelatihan dan akurasi melalui kombinasi \textit{Fused-MBConv}
dan strategi \textit{progressive learning} \cite{tan2021efficientnetv2}.
Penelitian pada domain klasifikasi citra bencana menunjukkan bahwa
keluarga \textit{EfficientNet} secara konsisten mengungguli berbagai
arsitektur CNN lainnya \cite{alam2021robust}. Evolusi CNN berlanjut
dengan hadirnya \textit{ConvNeXt} yang mengadopsi prinsip-prinsip
desain dari arsitektur \textit{Transformer} ke dalam kerangka CNN,
menghasilkan model yang lebih kuat sekaligus mempertahankan sifat
konvolusional \cite{liu2022convnext}. \par

Di sisi lain, \textit{Vision Transformer} (ViT) memperkenalkan
paradigma baru dengan menerapkan mekanisme \textit{self-attention} pada
sekuens \textit{patch} citra, memungkinkan penangkapan dependensi jarak
jauh sejak \textit{layer} pertama tanpa dibatasi oleh \textit{receptive
field} seperti pada CNN \cite{dosovitskiy2021image}. Perkembangan lebih
lanjut menghasilkan \textit{Swin Transformer} yang mengintegrasikan
konsep hierarki dan lokalitas dari CNN ke dalam arsitektur
\textit{Transformer} melalui mekanisme \textit{shifted windows},
menghasilkan representasi fitur yang lebih kaya dengan kompleksitas
komputasi yang lebih terkendali \cite{liu2021swin}. Dengan demikian,
keempat arsitektur yang digunakan dalam penelitian ini (yaitu
\textit{EfficientNetV2-M}, \textit{ConvNeXt-Base}, \textit{Swin-Small},
dan ViT-B/16) mewakili ragam arsitektur dari CNN murni hingga
\textit{Transformer} murni dengan karakteristik yang berbeda dalam
menangkap pola lokal maupun konteks global. Penelitian terbaru
menunjukkan bahwa arsitektur berbasis \textit{Transformer} seperti
\textit{CrisisViT} mulai diterapkan pada domain respons bencana dengan
performa yang kompetitif \cite{long2023crisisvit}, namun studi
komparatif yang mencakup ragam arsitektur CNN hingga \textit{Transformer}
secara sistematis pada \textit{dataset} \textit{CrisisMMD} belum pernah
dilakukan. \par

Mengingat kekuatan yang saling melengkapi dari paradigma CNN dan
\textit{Transformer}, \textit{ensemble learning} menjadi pendekatan yang
menjanjikan untuk memaksimalkan performa klasifikasi. \textit{Ensemble
learning} menggabungkan prediksi dari beberapa model untuk menghasilkan
prediksi yang lebih \textit{robust} dan akurat dibandingkan model
tunggal manapun \cite{rokach2010ensemble}. Penelitian terdahulu pada
\textit{CrisisMMD} menunjukkan bahwa \textit{ensemble} dari beberapa
model CNN dapat memberikan peningkatan akurasi dibandingkan model
tunggal \cite{firmansyah2022ensemble}, namun pendekatan tersebut
terbatas pada \textit{homogeneous ensemble} dari arsitektur yang
sejenis. \textit{Heterogeneous ensemble} yang menggabungkan model CNN
dan \textit{Transformer} berpotensi memanfaatkan kekuatan dari kedua
paradigma secara bersamaan, mengingat keragaman arsitektur yang lebih
tinggi cenderung menghasilkan peningkatan performa yang lebih signifikan
\cite{ganaie2022ensemble}, namun potensi ini belum dieksplorasi secara
mendalam untuk domain klasifikasi informasi krisis. \par

Selain pertanyaan mengenai konfigurasi model mana yang terbaik,
terdapat kebutuhan untuk memahami faktor-faktor yang berkontribusi
terhadap performa tersebut. Strategi \textit{fine-tuning} yang
digunakan dan komposisi model dalam \textit{ensemble} merupakan faktor
yang dapat mempengaruhi hasil secara signifikan, namun analisis sistematisnya dalam bentuk \textit{ablation study} 
masih jarang dilakukan pada domain klasifikasi citra bencana.
Pemahaman mendalam tentang kontribusi masing-masing faktor ini penting
agar penelitian tidak hanya menghasilkan angka performa tertinggi,
tetapi juga memberikan wawasan yang dapat direplikasi dan dikembangkan
lebih lanjut oleh komunitas peneliti. \par

Berdasarkan uraian di atas, penelitian ini melakukan analisis performa
secara komprehensif antara \textit{ensemble} heterogen dan model tunggal
untuk klasifikasi informasi krisis pada \textit{dataset}
\textit{CrisisMMD}. Empat arsitektur yang mewakili ragam CNN hingga
\textit{Transformer} (\textit{EfficientNetV2-M}, \textit{ConvNeXt-Base},
\textit{Swin-Small}, dan ViT-B/16) dievaluasi secara individual maupun
dalam berbagai konfigurasi \textit{ensemble} pada ketiga tugas
klasifikasi \textit{CrisisMMD}. Penelitian ini juga melakukan
\textit{ablation study} untuk menganalisis faktor-faktor yang
mempengaruhi performa, sehingga memberikan kontribusi yang tidak hanya
berupa hasil komparasi, tetapi juga pemahaman yang lebih dalam tentang
perilaku \textit{ensemble} heterogen di domain respons bencana. \par


\section{Rumusan Masalah} \label{I.RumusanMasalah}
Berdasarkan latar belakang yang telah diuraikan, rumusan masalah dalam
penelitian ini adalah: \par

\begin{enumerate}[noitemsep]
    \item Bagaimana performa masing-masing model tunggal yang mewakili
    ragam arsitektur CNN hingga \textit{Transformer}
    (\textit{EfficientNetV2-M}, \textit{ConvNeXt-Base},
    \textit{Swin-Small}, dan ViT-B/16) dalam klasifikasi informasi
    krisis pada ketiga tugas \textit{CrisisMMD}?

    \item Apakah \textit{ensemble} heterogen yang menggabungkan model
    CNN dan \textit{Transformer} mampu mengungguli performa model
    tunggal terbaik maupun \textit{ensemble} homogen, dan strategi
    \textit{ensemble} mana yang paling efektif untuk klasifikasi
    informasi krisis?

    \item Faktor-faktor apa saja yang secara signifikan mempengaruhi
    performa \textit{ensemble} heterogen, dan bagaimana kontribusi
    masing-masing faktor tersebut terhadap hasil klasifikasi
    berdasarkan \textit{ablation study}?
\end{enumerate}


\section{Tujuan Penelitian} \label{I.Tujuan}
Berdasarkan rumusan masalah yang telah diuraikan, tujuan dari penelitian
ini adalah: \par

\begin{enumerate}[noitemsep]
    \item Mengimplementasikan dan mengevaluasi empat arsitektur model
    tunggal (\textit{EfficientNetV2-M}, \textit{ConvNeXt-Base},
    \textit{Swin-Small}, dan ViT-B/16) pada ketiga tugas klasifikasi
    \textit{CrisisMMD} untuk menganalisis performa masing-masing
    arsitektur dalam mewakili ragam CNN hingga \textit{Transformer}.

    \item Mengimplementasikan dan mengevaluasi berbagai konfigurasi
    \textit{ensemble} heterogen dengan membandingkan performa
    \textit{ensemble} terhadap model tunggal terbaik dan
    \textit{ensemble} homogen, menggunakan strategi \textit{simple
    averaging}, \textit{weighted voting}, dan \textit{stacking}.

    \item Melakukan \textit{ablation study} untuk mengidentifikasi dan
    mengukur kontribusi faktor-faktor yang mempengaruhi performa
    \textit{ensemble} heterogen, sehingga menghasilkan pemahaman yang
    dapat direplikasi dan dikembangkan lebih lanjut.
\end{enumerate}


\section{Batasan Masalah} \label{I.Batasan}
Agar penelitian ini lebih fokus dan terarah, ditetapkan batasan masalah
sebagai berikut: \par

\begin{enumerate}[noitemsep]
    \item Penelitian hanya berfokus pada klasifikasi citra dari media
    sosial untuk respons bencana, tidak termasuk analisis teks atau
    pendekatan multimodal.

    \item Arsitektur model yang digunakan dibatasi pada empat model yang
    mewakili ragam CNN hingga \textit{Transformer}, yaitu
    \textit{EfficientNetV2-M}, \textit{ConvNeXt-Base},
    \textit{Swin-Small}, dan ViT-B/16.

    \item Strategi \textit{ensemble} yang dievaluasi dibatasi pada
    \textit{simple averaging}, \textit{weighted voting}, dan
    \textit{stacking} dengan \textit{Logistic Regression} sebagai
    \textit{meta-learner}.

    \item \textit{Dataset} yang digunakan adalah \textit{CrisisMMD}
    yang mencakup tujuh jenis bencana alam pada tiga tugas klasifikasi:
    \textit{Informative Classification} (Task 1), \textit{Humanitarian
    Categories} (Task 2), dan \textit{Damage Severity Assessment}
    (Task 3).

    \item \textit{Ablation study} dilakukan untuk menganalisis
    faktor-faktor yang mempengaruhi performa model, dengan cakupan
    aspek yang ditentukan berdasarkan hasil eksperimen awal.

    \item Analisis efisiensi komputasi (\textit{inference time},
    \textit{memory usage}, FLOPs) dilakukan sebagai analisis pendukung
    pada Task 1, bukan sebagai tujuan utama penelitian.

    \item Penelitian tidak mencakup penerapan sistem ke lingkungan
    produksi atau pengembangan aplikasi pengguna akhir.
\end{enumerate}


\section{Manfaat Penelitian} \label{I.Manfaat}
Manfaat yang diharapkan dari penelitian ini adalah: \par

\subsection{Manfaat Teoritis}
\begin{itemize}[noitemsep]
    \item Memberikan kontribusi berupa studi komparatif sistematis yang
    mencakup ragam arsitektur CNN hingga \textit{Transformer} pada
    domain klasifikasi informasi krisis, yang selama ini belum pernah
    dilakukan secara komprehensif pada \textit{dataset}
    \textit{CrisisMMD}.

    \item Memperkaya literatur mengenai \textit{ensemble learning}
    dengan mengeksplorasi efektivitas \textit{ensemble} heterogen yang
    menggabungkan paradigma CNN dan \textit{Transformer}, serta
    membandingkannya secara langsung dengan \textit{ensemble} homogen.

    \item Menyediakan hasil \textit{ablation study} yang dapat menjadi
    referensi bagi penelitian selanjutnya dalam memahami faktor-faktor
    yang menentukan keberhasilan \textit{ensemble} heterogen di domain
    respons bencana.
\end{itemize}

\subsection{Manfaat Praktis}
\begin{itemize}[noitemsep]
    \item Memberikan panduan berbasis bukti eksperimen bagi praktisi
    dalam memilih konfigurasi model yang sesuai untuk sistem
    klasifikasi informasi krisis, dengan mempertimbangkan karakteristik
    masing-masing arsitektur dan kebutuhan penerapan.

    \item Membantu pengembangan sistem respons bencana berbasis
    \textit{Artificial Intelligence} (AI) yang lebih akurat melalui
    pemahaman yang lebih dalam tentang kekuatan yang saling melengkapi
    antara arsitektur CNN dan \textit{Transformer} dalam konteks
    \textit{ensemble}.

    \item Menjadi referensi implementasi sistem klasifikasi otomatis
    yang dapat diadaptasi untuk berbagai skenario klasifikasi citra
    bencana dengan karakteristik yang serupa.
\end{itemize}


\section{Sistematika Penulisan} \label{I.Sistematika}
Sistematika penulisan tugas akhir ini disusun sebagai berikut: \par

\subsection*{Bab I}
Bab ini membahas latar belakang penelitian, rumusan masalah, tujuan
penelitian, batasan masalah, manfaat penelitian, dan sistematika
penulisan. \par

\subsection*{Bab II}
Bab ini menguraikan tinjauan pustaka dari penelitian terdahulu serta
dasar teori yang menjadi landasan penelitian, termasuk klasifikasi
citra untuk respons bencana, \textit{dataset} \textit{CrisisMMD},
arsitektur \textit{EfficientNetV2}, \textit{ConvNeXt}, \textit{Swin
Transformer}, \textit{Vision Transformer}, \textit{ensemble learning},
\textit{transfer learning}, serta metrik evaluasi. \par

\subsection*{Bab III}
Bab ini menjelaskan alur penelitian, penjabaran langkah-langkah
penelitian, alat dan bahan yang digunakan, metode pengembangan yang
meliputi preparasi \textit{dataset}, konfigurasi model, strategi
\textit{ensemble}, rancangan \textit{ablation study}, serta rancangan
pengujian. \par

\subsection*{Bab IV}
Bab ini menyajikan hasil implementasi dan pengujian yang telah
dilakukan, analisis performa model tunggal, analisis performa
\textit{ensemble} dibandingkan model tunggal dan \textit{ensemble}
homogen, hasil \textit{ablation study}, analisis pendukung efisiensi
komputasi, serta pembahasan mendalam terhadap seluruh hasil yang
diperoleh. \par

\subsection*{Bab V}
Bab ini berisi kesimpulan dari hasil penelitian yang menjawab seluruh
rumusan masalah, serta saran untuk pengembangan penelitian
selanjutnya. \par
